% ============================ REMERCIEMENTS ==========================
\chapter*{Remerciements}
\addcontentsline{toc}{chapter}{Remerciements}

Au terme de ce stage de Projet de Fin d'Année, je tiens à exprimer ma profonde gratitude à
l'ensemble des personnes qui ont contribué, de près ou de loin, à son bon déroulement.

Mes remerciements s'adressent en premier lieu à M. DARKAOUI Hatim, mon encadrant au sein
d'ALTEN, pour la confiance qu'il m'a accordée, la pertinence de ses orientations techniques
et sa disponibilité tout au long de ces huit semaines. Ses retours réguliers m'ont permis de
garder un cap clair entre l'ambition fonctionnelle du projet et les contraintes réelles d'un
environnement industriel.

Je remercie également l'ensemble des équipes d'ALTEN pour leur accueil chaleureux et pour
l'environnement de travail stimulant dans lequel j'ai pu évoluer. Travailler au contact
d'ingénieurs confrontés quotidiennement à des problématiques concrètes de gestion de
documentation technique a donné tout son sens à ce projet.

Ma reconnaissance va aussi au corps professoral de l'École Nationale Supérieure d'Arts et
Métiers de Meknès, et en particulier aux enseignants de la filière Intelligence Artificielle
\& Technologies des Données, dont les enseignements ont constitué le socle théorique sur
lequel ce travail s'est appuyé.

Enfin, je remercie ma famille et mes proches pour leur soutien constant.

% ============================ RÉSUMÉ =================================
\cleardoublepage
\chapter*{Résumé}
\addcontentsline{toc}{chapter}{Résumé}

Ce rapport présente la conception et la réalisation d'un \textbf{assistant de recherche
documentaire multimodal} fondé sur le paradigme RAG (\emph{Retrieval-Augmented Generation}),
développé au cours d'un stage PFA de huit semaines au sein d'ALTEN. L'objectif est de
permettre à un ingénieur d'interroger en langage naturel un corpus de documents techniques
(articles scientifiques, manuels d'ingénierie au format PDF ou Word) et d'obtenir des
réponses \emph{sourcées}, combinant extraits textuels cités et figures pertinentes affichées
directement dans l'interface.

La chaîne de traitement réalisée couvre l'intégralité du cycle de vie documentaire :
analyse structurelle des documents avec \textbf{Docling} (détection de mise en page,
extraction des figures, reconstruction des tableaux), découpage sémantique par
\emph{HybridChunker}, description automatique des figures par un modèle de vision-langage
(\textbf{Gemma-3 27B}), vectorisation unifiée des chunks textuels et des descriptions
d'images par \textbf{mistral-embed}, puis indexation dans \textbf{ChromaDB}. La phase de
restitution repose sur une \textbf{recherche hybride} combinant similarité dense et BM25 par
fusion de rangs réciproques (RRF), un routage automatique de la question vers les groupes
documentaires pertinents par LLM, et la génération de la réponse par
\textbf{mistral-small-latest} avec citation systématique des sources. Les capacités sont
exposées par une API REST \textbf{FastAPI} et une interface \textbf{React/TypeScript}
sécurisée par \textbf{Clerk}.

La qualité du système a été évaluée avec le cadre \textbf{RAGAS}, en utilisant un juge LLM
indépendant (Llama~3.3~70B). Les résultats obtenus (\emph{context recall} de 1{,}00,
\emph{answer relevancy} de 0{,}93 et \emph{faithfulness} moyenne de 0{,}83) valident
l'efficacité du retrieval hybride et de l'indexation multimodale, tout en identifiant des
axes d'amélioration précis.

\vspace{3mm}
\noindent\textbf{Mots-clés :} RAG multimodal, LLM, embeddings, recherche hybride, BM25,
RRF, Docling, ChromaDB, Mistral AI, VLM, FastAPI, React, RAGAS.

% ============================ ABSTRACT ===============================
\cleardoublepage
\chapter*{Abstract}
\addcontentsline{toc}{chapter}{Abstract}

This report presents the design and implementation of a \textbf{multimodal
Retrieval-Augmented Generation (RAG) assistant}, developed during an eight-week internship
at ALTEN. The system enables engineers to query a corpus of technical documents (research
papers and engineering manuals in PDF or Word format) in natural language, and to receive
\emph{grounded} answers combining cited text excerpts and the relevant figures displayed
directly in the user interface.

The implemented pipeline covers the full document lifecycle: structural parsing with
\textbf{Docling} (layout detection, figure extraction, table reconstruction), semantic
chunking with the \emph{HybridChunker}, automatic figure description using a
vision-language model (\textbf{Gemma-3 27B}), unified vectorization of text chunks and
image descriptions with \textbf{mistral-embed}, and indexing into \textbf{ChromaDB}.
Retrieval combines dense similarity search and BM25 keyword matching through reciprocal
rank fusion (RRF), augmented by automatic LLM-based routing of each question to the
relevant document groups. Answers are generated by \textbf{mistral-small-latest} with
systematic source citation. All capabilities are exposed through a \textbf{FastAPI} REST
API and a \textbf{React/TypeScript} interface secured with \textbf{Clerk} authentication.

System quality was evaluated with the \textbf{RAGAS} framework, using an independent LLM
judge (Llama~3.3~70B). The results (context recall of 1.00, answer relevancy of 0.93 and
average faithfulness of 0.83) validate the effectiveness of hybrid retrieval and
multimodal indexing, while identifying precise improvement paths.

\vspace{3mm}
\noindent\textbf{Keywords:} multimodal RAG, LLM, embeddings, hybrid search, BM25, RRF,
Docling, ChromaDB, Mistral AI, VLM, FastAPI, React, RAGAS.
